% =========================
% MAU Master's Thesis: Abdalazeez Asaad
% Rejoinder to examiner and opponent comments
% =========================
\documentclass[11pt,a4paper]{article}

% --- Page / typography ---
\usepackage[a4paper,margin=30mm]{geometry}
\usepackage[utf8]{inputenc}
\usepackage[T1]{fontenc}
\usepackage{lmodern}
\usepackage{setspace}
\setstretch{1.0}
\usepackage{amsmath}
\usepackage{microtype}
\usepackage[hidelinks]{hyperref}
\usepackage{booktabs}
\usepackage{longtable}
\usepackage{array}
\usepackage{ragged2e}
\usepackage{enumitem}

\pagestyle{plain}
\justifying

\newcommand{\rjcomment}[1]{\par\noindent\textbf{Comment:}\ #1\par}
\newcommand{\rjresponse}[1]{\par\noindent\textbf{Response: #1}\par}

\title{Rejoinder}
\author{Abdalazeez Asaad}
\date{}

\begin{document}

\maketitle

Responses to comments received on the thesis draft from the examiner (Erdal Akin) and the
opponents (linne, Pari). Organised in two parts, in the order the comments appear in the
thesis. Each entry states the comment, whether it was fixed or is defended as written, and
what actually changed. References are to section names, not line numbers, since the
authoritative working copy is maintained outside this repository.

% =====================================================
\section{Part 1: Examiner Comments}
% =====================================================

\subsection{E1: Section 2.5 (Research Gaps and Positioning)}

\rjcomment{Phrase cautiously if you cannot support by a systematic review, as you have not
done for this work.}

\rjresponse{Fixed.}
The claim is now scoped to the literature reviewed in this thesis (not an absolute claim about
the field), and specifically extended to note that no distance-accuracy analysis exists for the
Swedish SMHI network either, citing Segersson et al.\ and the Stockholm/Gothenburg/Malm\"o
dispersion study to show the Swedish literature was checked.

\subsection{E2: Section 3.3 (Spatial Validation Protocol)}

\rjcomment{Using buffered spatial leave-one-out validation is appropriate, but the 5\,km buffer
is much smaller than the reported spatial autocorrelation range (150--200\,km). Therefore, the
validation should not be described as fully ``spatially honest.'' The limitation should be
stated more clearly, and a small sensitivity test with different feasible buffer distances would
strengthen the work.}

\rjresponse{Fixed.}
Three changes:
\begin{enumerate}[label=\arabic*.,leftmargin=*]
  \item The phrase ``spatially honest'' is removed throughout (Introduction and \S3.3) and
    replaced with language that does not overstate the protocol (e.g.\ ``with near-neighbour
    leakage suppressed'').
  \item The existing buffer-justification paragraph in \S3.3 is sharpened to state explicitly
    that the 5\,km buffer reduces but does not eliminate the optimistic bias from spatial
    autocorrelation, and that the reported errors should be read as a lower bound on true
    prediction error.
  \item A buffer-distance sensitivity analysis was run: the full buffered SLOO procedure was
    repeated at 0, 5, 10, 25, and 50\,km exclusion radii
    (\texttt{data/phase6\_buffer\_sensitivity.py}). The estimator ranking is unchanged at every
    buffer tested (IDW best for PM$_{2.5}$, Random Forest best for NO$_2$, LUR worst for both),
    and mean RMSE for the two selected estimators varies by no more than 6\,\% between the 5\,km
    and 50\,km buffers. A new table and paragraph reporting this are added to \S3.3, and a
    corresponding sentence is added to the limitations section (\S7.4). Buffers approaching the
    150--200\,km autocorrelation range are not testable: above roughly 80\,km, at least one
    validation fold loses all training data, which is now stated explicitly.
\end{enumerate}

\subsection{E3: Section 4.1 (Air Quality Observations)}

\rjcomment{This is a small effective spatial sample for generalization.}

\rjresponse{Defended, no change.}
Agreed, and already stated as a limitation. \S7.4 notes that the 11 concurrently-passing
stations are ``substantially smaller than the 20--40 monitoring sites comprising each ESCAPE
study area,'' constraining the statistical power of the validation and the granularity of the
resulting decay curves. No further change made; this is treated as an acknowledged constraint of
the Swedish national network rather than a defect to correct.

\subsection{E4: Section 5.1 (NO2 Estimation Results, Table 5)}

\rjcomment{In table 5, IDW provides better RMSE in 7/11 stations. Even RF provides lower RMSE,
this situation needs to be clarified. It would be great if you could provide $R^2$ and MAE of
LUR and IDW.}

\rjresponse{Fixed.}
Added MAE and $R^2$ for LUR and IDW (previously only RMSE was shown). Added a short explanation
of why IDW wins at 7/11 stations despite RF's lower mean RMSE: RF's mean is pulled down by two
isolated rural stations where it performs exceptionally well, while IDW is competitive or better
at most of the remaining nine.

\subsection{E5: Section 5.2 (PM2.5 Estimation Results, Table 6)}

\rjcomment{As seen in Table 6 also, IDW provides better performance for all the aspects. So, it
would be better to provide all the related metrics for the three algorithms and compare fairly
instead of making RF the core algorithm.}

\rjresponse{Fixed.}
Added MAE and $R^2$ for LUR and IDW here as well. Noted that no equivalent caveat is needed for
PM$_{2.5}$: IDW wins at 10/11 stations, which is consistent with its lower mean RMSE (unlike the
NO$_2$ case).

\subsection{E6: Section 6.2 (Decay Curve Fitting)}

\rjcomment{These criteria should be defined and elaborated to show the effect and necessity.
Also, this section needs a transition why we need to see Decay curve fitting.}

\rjresponse{Fixed.}
AIC/BIC are now defined and elaborated in \S6.2 (formulas, what they measure, why AIC over raw
fit), with the design rationale for the choice moved to \S3.4 so it is argued once, in the
methodology, rather than repeated. Added a transition at the start of \S6.2 linking it back to
\S3.4.

\subsection{E7: Section 6.4 (Reliable Prediction Distance Threshold and Answer to RQ2)}

\rjcomment{I would therefore recommend that 64\,km should not be called ``the reliable
prediction distance'' without qualification. It would be better described as something like: the
estimated prediction-distance threshold under the adopted 50\,\% RMSE criterion. This
distinction should appear in the Abstract, Results, Discussion and Conclusion. The 64\,km and
especially the 6\,km values should therefore be presented as criterion-dependent estimates, not
precise national reliability limits.}

\rjresponse{Fixed.}
The thesis presents 64\,km and 6\,km as ``estimated prediction-distance thresholds under the
adopted 50\,\% criterion,'' not fixed reliability limits, with an added note that both are
criterion-dependent estimates (with extra caution on the 6\,km NO$_2$ figure). This is applied
in the Abstract, in \S6.4 (retitled ``Estimated Prediction-Distance Threshold and Answer to
RQ2''), in the Discussion's RQ2-answer paragraph, and in the Conclusion.

On review, three further occurrences of the unqualified phrase ``the reliable prediction
distance/threshold'' attached directly to the 64\,km figure had survived outside \S6.4, in the
Placement section's opening paragraph, in the Integrated Results (RQ1--RQ3) paragraph (two
occurrences), and in the Discussion's Risks and Constraints subsection, the last of which also
asserted that a covered location was ``trustworthy'' without qualification. All three are now
corrected to ``estimated prediction-distance threshold,'' and the Risks and Constraints sentence
is reworded to state that coverage indicates the location is expected to meet the adopted 50\,\%
criterion, not that it is unqualifiedly ``trustworthy.''

\subsection{E8: Section 7.1 (National Coverage Surface)}

\rjcomment{Good point to declare. \emph{(on the statement that NO2 baseline coverage is minimal:
7 of 4,588 land cells, 0.2\%)}}

\rjresponse{Acknowledged, no change requested.}
This was agreement with a point already made in the thesis, not a request for revision.

\subsection{E9: Section 7.3 (Coverage Evaluation and DSR Assessment)}

\rjcomment{Can you extend the work for population-weighted coverage. You can use same experiment
you have done for this evaluation and do again wrt population. I would like to see this. If time
permits, you can also work on urban-area coverage, which is pretty similar to the
population-based approach.}

\rjresponse{Fixed.}
Urban-area coverage was already reported in the thesis (land coverage 10.9\,\% $\rightarrow$
67.2\,\%; urban coverage 38.8\,\% $\rightarrow$ 71.4\,\%). Population-weighted coverage is new:
each 10\,km grid cell was weighted by its resident population, aggregated from the JRC GEOSTAT
2018 1\,km population grid ($\sim$10.6 million inhabitants assigned across the national land
grid; \texttt{data/phase7\_population\_coverage.py}). Results: population-weighted PM$_{2.5}$
coverage rises from 45.0\,\% to 77.1\,\% nationally (67.1\,\% to 75.9\,\% within urban fabric)
after the 20 recommended placements, markedly higher than the land-area figures, because
existing stations are concentrated in populated areas. For NO$_2$, population-weighted coverage
moves only from 10.1\,\% to 10.4\,\% nationally and is unchanged at 17.0\,\% within urban fabric,
consistent with the land-area result: the PM$_{2.5}$-optimised deployment does not materially
improve NO$_2$ coverage on any measure. A new paragraph and table (land-area vs.\
population-weighted, national vs.\ urban, before/after, both pollutants) are added to \S7.3, and
a caveat is added to \S7.4 noting that the population-weighted NO$_2$ figures inherit the same
10\,km grid-resolution limitation already stated for the land-area NO$_2$ figures.

\subsection{E10: Figures 2 and 3, decay curves (LUR centre panel)}

\rjcomment{Explain more about the middle (LUR) panel: the fitted line is nearly straight/flat,
which could be read as LUR being more stable with respect to distance. Why was LUR not
selected?}

\rjresponse{Fixed.}
The single existing sentence in \S6.3 (``LUR curves are reported for completeness but carry
substantial uncertainty due to the high cross-fold variability of the linear model'') is
expanded to explain the flat line explicitly. The LUR fit is flat because its error is uniformly
poor, not because it is distance-robust: LUR RMSE sits at approximately 15\,$\mu$g/m$^3$
(NO$_2$) and 6\,$\mu$g/m$^3$ (PM$_{2.5}$) at every distance, in both cases above the pollutant's
own mean concentration (12.85 and $\sim$5.3\,$\mu$g/m$^3$). The near-zero slope arises because
LUR's error is dominated by cross-fold instability (30 partly collinear predictors on 11 folds),
which is large relative to any distance effect; the wide vertical scatter of the points and the
poor curve fit (AIC 83.5 / 58.8 versus 60.7--63.4 and 25.3--38.6 for RF and IDW) show the flat
line is not an informative relationship. LUR is therefore not a candidate for the threshold,
consistent with it being the least accurate estimator for both pollutants in \S5. A one-clause
note to the same effect is added to the captions of Figures 2 and 3.

% =====================================================
\section{Part 2: Opponent Comments}
% =====================================================

\subsection{O1: Title page}

\rjcomment{\emph{(linne)} Shouldn't this be a sub-title? \emph{(on the portion of the title
after the colon)}}

\rjresponse{Defended, no change.}
The title is the one approved in the research proposal; the wording is unchanged either way.
Reformatting it as a two-line title/sub-title is a cosmetic option, not a required change, and is
left as a single line consistent with the approved proposal.

\subsection{O2: Section 2 (Background and Related Work)}

\rjcomment{\emph{(linne)} A short introduction to the different topics you talk about in this
section would be nice.}

\rjresponse{Fixed.}
A short paragraph was added immediately after the section heading, previewing the five
subsections that follow (regulatory monitoring, low-cost IoT sensors, spatial estimation
methods, sensor placement optimisation, and the research gaps).

\subsection{O3: Section 2 / 2.1 (Background and Related Work)}

\rjcomment{\emph{(Pari)} The section discusses relevant previous studies but the process used to
identify and select the literature is not described. Could you briefly explain which databases or
search tools were used.}

\rjresponse{Fixed.}
A paragraph was added describing the search process: literature was identified primarily through
Google Scholar, with IEEE Xplore used for the conference proceedings, and by following the
reference lists of the principal review papers (Hoek et al.; the ESCAPE studies) to trace
related work. Search terms are listed. The paragraph states explicitly that this was a targeted
narrative review, not a systematic review, consistent with the hedge already applied to \S2.5
(E1 above).

\subsection{O4: Section 2.1 (Urban Air Quality Monitoring and Regulatory Networks)}

\rjcomment{\emph{(linne)} This thesis? \emph{(on the phrase ``The present thesis'')}}

\rjresponse{Fixed.}
``The present thesis'' is reworded to ``This thesis'' at this location and at one further,
equivalent occurrence in \S2.3.

\subsection{O5: Section 2.3 (Spatial Estimation of Air Quality at Unmonitored Locations)}

\rjcomment{\emph{(linne)} When writing abbreviations the text should be ``Inverse Distance
Weighting.'' By some university writing guidelines for acronyms and abbreviations.}

\rjresponse{Defended, no change.}
The term is introduced in the standard form used throughout the thesis: lower case at first use
with the acronym given in parentheses, and capitalised in the List of Abbreviations. This
follows the convention used consistently for every other abbreviation in the thesis (e.g.\
``land-use regression (LUR)''), so it is retained for consistency rather than treated as an
isolated exception.

\subsection{O6: Section 2.3 (Spatial Estimation of Air Quality at Unmonitored Locations)}

\rjcomment{\emph{(linne)} Reference? \emph{(on the paragraph describing LUR's known
limitations)}}

\rjresponse{Fixed.}
The paragraph is attributed to Hoek et al.\ and reworded to state only what that review
supports: that LUR models are empirical and area-specific, and their predictive performance is
not guaranteed to transfer to areas or periods outside the monitoring campaign used to fit them,
alongside the existing point that a linear form does not represent non-linear predictor
interactions by construction. The citation metadata was confirmed against ADS and ScienceDirect,
and the transferability limitation is consistently attributed to this review in the secondary
literature, so no verification marker remains in the thesis source.

\subsection{O7: Section 3 (Research Methodology)}

\rjcomment{\emph{(linne)} All theses do not need to look exactly the same, but I think
best-practice is that new ``chapters'' should start on a new page.}

\rjresponse{Defended, no change.}
The opponent's own phrasing acknowledges this is a formatting preference, not a requirement. The
document follows the MAU thesis template's layout conventions. Left as is; a one-line
\texttt{\textbackslash clearpage} change is a trivial future option if a stronger visual break is
wanted.

\subsection{O8: Section 3.2.2 (Feature Set)}

\rjcomment{\emph{(Pari)} Section 4.1 represents that all 34 stations may contribute training
data, while only 11 stations are eligible as validation targets. So, the phrase ``11 unique
spatial training locations'' is unclear and appears inconsistent with Sections 4.1 and 5.}

\rjresponse{Fixed.}
This was a genuine wording inconsistency; \S4.1 already states the correct relationship (34
stations may contribute training rows; 11 concurrently-passing stations are eligible to be
withheld as validation targets), and \S5 confirms training uses 26--33 stations per fold, not
11. The phrase ``11 unique spatial training locations'' is corrected to ``11 spatially
independent validation targets'' (or equivalent) at the three locations where it previously
appeared incorrectly (\S3.2.2, \S3.4, \S7.4), aligning them with \S4.1.

\subsection{O9: Section 5 (Spatial Estimation Models)}

\rjcomment{\emph{(linne)} What is ``this chapter''? I think the structure of what is results
could perhaps be better structured or perhaps just named differently so it is more clear what is
the actual results of the study. I cannot see that chapter is used anywhere else in the thesis
and it is therefore unclear what is chapter/section.}

\rjresponse{Partially fixed.}
``This chapter'' was a genuine slip: it is the only use of ``chapter'' in the entire thesis,
which otherwise uses ``section'' as the top-level unit throughout. It is corrected to ``the
results'' with an explicit section range. The broader structural point is defended: the three
results sections are each named by research question (RQ1/RQ2/RQ3) and explicitly signposted in
the opening paragraph of \S5, which is judged sufficient without renaming the sections
themselves.

\subsection{O10: Section 5 (Spatial Estimation Models)}

\rjcomment{\emph{(linne)} This is ``double negative'' phrasing which is a bit unclear, and as
per my understanding should be avoided. \emph{(on ``not excluded by the 5\,km buffer'')}}

\rjresponse{Fixed.}
Reworded to ``the stations that remained after the 5\,km exclusion buffer was applied,''
removing the double negative.

\subsection{O11: Section 6.4 (Reliable Prediction Distance Threshold and Answer to RQ2)}

\rjcomment{\emph{(Pari)} The sensitivity analysis clearly explains that the 50\,\% criterion is
a modelling choice. But, it is still unclear why this level represents the ``minimum
requirement'' for spatial planning.}

\rjresponse{Fixed.}
The unsupported claim that 50\,\% is the ``minimum requirement for spatial planning purposes''
is removed. The paragraph is reworded to state plainly that the 50\,\% level is a working
threshold adopted for this study, not a universal planning standard, and to point immediately to
the sensitivity analysis that tests it, consistent with the criterion-dependent framing already
applied elsewhere in \S6.4 (E7 above).

\subsection{O12: Figure 6}

\rjcomment{\emph{(linne)} The triangles are red not green. A little confusion to write ``placed
sensors'' on the figure, and ``recommended sensor locations'' in the description? There are no
red circles in the figure.}

\rjresponse{Fixed.}
The figure itself (\texttt{fig\_coverage\_before\_after.png}) already used the correct, current
colour scheme (blue circles for existing stations, red triangles for recommended sensors),
matching Figures 7 and 8; only the caption was stale, describing an earlier colour scheme (red
circles, green triangles). The caption is corrected to match the actual figure, and the ``placed
sensors'' legend label vs.\ ``recommended sensor locations'' caption wording is reconciled by
noting the legend label explicitly in the caption.

\subsection{O13: Figure 7}

\rjcomment{\emph{(linne)} These labels are impossible to see or read. \emph{(on the numbered
placement-rank labels)}}

\rjresponse{Fixed.}
Confirmed: the rank labels were drawn at font size 7 in the same colour as the marker, directly
on top of the darkest part of the heatmap, and were effectively illegible. The figure
(\texttt{fig\_placement\_map.png}) is regenerated with larger, bold, black labels with a white
halo, clearly legible against both the heatmap and the markers. A new standalone script,
\texttt{data/phase7\_fix\_placement\_labels.py}, rebuilds this figure from saved results without
touching any other figure or data file.

\subsection{O14: References}

\rjcomment{\emph{(Pari)} IEEE reference guide is to write et al.\ after six authors.}

\rjresponse{Fixed.}
Four bibliography entries exceeded six authors (Hoek et al.\ 2008, 7 authors; Segersson et al.\
2017, 7 authors; Roberts et al.\ 2017, 14 authors; Kilbo Edlund et al.\ 2024, 14 authors). Each
is now truncated to the first author plus ``and others'' in \texttt{references.bib}, which the
thesis's \texttt{ieeetr} bibliography style renders as ``et al.'' No other references or the
bibliography style itself were changed.

% =====================================================
\section{Summary}
% =====================================================

\begin{longtable}{@{}c p{0.52\textwidth} p{0.30\textwidth}@{}}
\toprule
\textbf{\#} & \textbf{Comment} & \textbf{Outcome} \\
\midrule
\endfirsthead
\toprule
\textbf{\#} & \textbf{Comment} & \textbf{Outcome} \\
\midrule
\endhead
\bottomrule
\endfoot
E1  & \S2.5 systematic-review hedging & Fixed \\
E2  & \S3.3 buffer / ``spatially honest'' & Fixed (wording + sensitivity table + limitation) \\
E3  & \S4.1 small validation sample & Defended \\
E4  & \S5.1 Table 5 IDW/RF & Fixed \\
E5  & \S5.2 Table 6 IDW/RF & Fixed \\
E6  & \S6.2 AIC/BIC definition & Fixed \\
E7  & \S6.4 threshold terminology & Fixed \\
E8  & \S7.1 NO2 baseline coverage & Acknowledged \\
E9  & \S7.3 population-weighted coverage & Fixed (new analysis) \\
E10 & Figures 2--3 LUR centre panel & Fixed (expanded \S6.3 + caption clauses) \\
O1  & Title sub-title & Defended \\
O2  & \S2 intro paragraph & Fixed \\
O3  & \S2/2.1 literature search process & Fixed \\
O4  & \S2.1 ``This thesis?'' wording & Fixed \\
O5  & \S2.3 IDW capitalisation & Defended \\
O6  & \S2.3 LUR limitations reference & Fixed \\
O7  & \S3 chapters on new page & Defended \\
O8  & \S3.2.2 ``11 unique spatial training locations'' & Fixed \\
O9  & \S5 ``this chapter'' / structure & Partially fixed (wording fixed; structure defended) \\
O10 & \S5 double negative & Fixed \\
O11 & \S6.4 ``minimum requirement'' & Fixed \\
O12 & Figure 6 caption & Fixed \\
O13 & Figure 7 labels & Fixed \\
O14 & References ``et al.'' & Fixed \\
\end{longtable}

\end{document}
